\documentclass[12pt,a4paper]{article}

% =========================================================
% PACKAGES
% =========================================================

\usepackage[a4paper,margin=1in]{geometry}
\usepackage{graphicx}
\usepackage{booktabs}
\usepackage{amsmath,amssymb}
\usepackage{float}
\usepackage{hyperref}
\usepackage{xcolor}
\usepackage{listings}
\usepackage{enumitem}
\usepackage{fancyhdr}
\usepackage{longtable}
\usepackage{array}

% =========================================================
% HYPERLINK SETTINGS
% =========================================================

\hypersetup{
    colorlinks=true,
    linkcolor=blue,
    urlcolor=blue,
    citecolor=blue
}

% =========================================================
% PAGE STYLE
% =========================================================

\pagestyle{fancy}
\fancyhf{}

\lhead{ICS1512}
\rhead{Machine Learning Algorithms Laboratory}
\cfoot{\thepage}

% =========================================================
% CODE LISTING SETTINGS
% =========================================================

\lstset{
    language=Python,
    basicstyle=\ttfamily\footnotesize,
    numbers=left,
    numberstyle=\tiny,
    frame=single,
    breaklines=true,
    showstringspaces=false,
    tabsize=4
}

% =========================================================
% DOCUMENT
% =========================================================

\begin{document}

% =========================================================
% TITLE / COVER INFORMATION
% =========================================================

\begin{center}

{\Large \textbf{MACHINE LEARNING ALGORITHMS LABORATORY}}\\[0.3cm]

{\large \textbf{Experiment 4}}\\[0.3cm]

{\Large \textbf{Binary Classification using Linear and Kernel-Based Models}}

\end{center}

\vspace{0.5cm}

\begin{tabular}{|p{4cm}|p{10.5cm}|}
\hline
\textbf{Experiment No.} & 4 \\ \hline

\textbf{Experiment Title} &
Binary Classification using Linear and Kernel-Based Models \\ \hline

\textbf{Student Name} & Harshitaa B\\ \hline

\textbf{Register Number} & 3122247001021 \\ \hline

\textbf{Faculty} & Dr. Poreddy Ajay Kumar Reddy \\ \hline

\textbf{Submission Date} & 16th August 2026 \\ \hline

\end{tabular}

\vspace{0.5cm}

\begin{center}
\textbf{Subject Code: ICS1512}\\
\textbf{Machine Learning Algorithms Laboratory}
\end{center}

\newpage

% =========================================================
% 1. OBJECTIVE
% =========================================================

\section{Objective}

The objectives of this experiment are:

\begin{enumerate}
    \item To implement binary classification using Logistic Regression
    and Support Vector Machine (SVM).
    
    \item To classify email samples into spam and ham classes using the
    Spambase dataset.
    
    \item To study the effect of regularization and other
    hyperparameters on Logistic Regression performance.
    
    \item To compare different SVM kernels including linear,
    polynomial, RBF and sigmoid kernels.
    
    \item To perform hyperparameter tuning using cross-validation.
    
    \item To evaluate the models using accuracy, precision, recall and
    F1-score.
    
    \item To compare Logistic Regression and SVM using five-fold
    cross-validation.
    
    \item To analyse the effect of model selection and hyperparameter
    tuning on classification performance.
\end{enumerate}

% =========================================================
% 2. PROBLEM STATEMENT
% =========================================================

\section{Problem Statement}

The objective of this experiment is to develop binary classification
models for identifying emails as either spam or ham. Logistic Regression
is used as a linear probabilistic classifier, while Support Vector
Machine is used as a margin-based classifier.

The performance of the models is evaluated before and after
hyperparameter tuning. Different SVM kernels are also compared to
understand their effect on classification performance.

The final models are evaluated using standard classification metrics
and five-fold cross-validation.

% =========================================================
% 3. DATASET DESCRIPTION
% =========================================================

\section{Dataset Description}

The Spambase dataset contains numerical features extracted from email
content along with a binary target indicating whether an email is spam
or non-spam (ham).

\begin{longtable}{|p{5cm}|p{9cm}|}
\hline
\textbf{Property} & \textbf{Description} \\ \hline

Dataset Name &
Spambase Dataset \\ \hline

Dataset Type &
Binary Classification \\ \hline

Application &
Spam and Ham Email Classification \\ \hline

Target Variable &
Spam / Ham \\ \hline

Feature Type &
Numerical features extracted from email content \\ \hline

Number of Classes &
2 \\ \hline

Dataset Source &
Spambase dataset \\ \hline

\end{longtable}

\textbf{Note:} The exact sample count, feature count and missing-value
count should be entered from the dataset inspection performed in the
notebook.

% =========================================================
% 4. EXPLORATORY DATA ANALYSIS
% =========================================================

\section{Exploratory Data Analysis}

Exploratory Data Analysis was performed before training the
classification models. The main purpose of EDA was to understand the
structure of the dataset and identify potential preprocessing
requirements.

The following aspects were considered:

\begin{itemize}
    \item Dataset dimensions
    \item Data types
    \item Missing values
    \item Statistical summary
    \item Class distribution
    \item Feature distributions
    \item Correlation between features
\end{itemize}

\subsection{Dataset Overview}

The dataset was inspected to determine the number of samples, features,
data types and target classes.

\textbf{Observation:}

The dataset consists of numerical email-related features and a binary
target variable representing spam and ham emails. Understanding the
dataset structure is important before applying classification
algorithms.

% ---------------------------------------------------------
% CLASS DISTRIBUTION
% ---------------------------------------------------------

\subsection{Class Distribution}

\begin{figure}[H]
\centering
\includegraphics[width=0.75\linewidth]{class_distribution.png}
\caption{Distribution of Spam and Ham Emails}
\label{fig:class_distribution}
\end{figure}

\textbf{Observation:}

The class distribution graph shows the number of spam and ham samples
in the dataset. The distribution helps determine whether the dataset
contains a significant class imbalance. This information is useful
when interpreting accuracy along with precision, recall and F1-score.

% ---------------------------------------------------------
% CORRELATION MATRIX
% ---------------------------------------------------------

\subsection{Correlation Matrix}

\begin{figure}[H]
\centering
\includegraphics[width=0.80\linewidth]{correlation_matrix.png}
\caption{Correlation Matrix of Dataset Features}
\label{fig:correlation}
\end{figure}

\textbf{Observation:}

The correlation matrix represents the pairwise relationship between
the numerical features. Strong positive or negative correlations can
indicate that some features contain related information. The matrix
also provides an overview of the relationships among the input
variables before model training.

% =========================================================
% 5. DATA PREPROCESSING
% =========================================================

\section{Data Preprocessing}

The following preprocessing steps were performed:

\begin{enumerate}
    \item Missing values were checked in the dataset.
    \item The target variable was separated from the input features.
    \item The data was divided into training and testing sets.
    \item Numerical features were standardized before model training.
    \item The same preprocessing procedure was applied consistently
    during model evaluation.
\end{enumerate}

Feature standardization is particularly important for Logistic
Regression and SVM because both models are affected by the scale of
the input features.

% =========================================================
% 6. MACHINE LEARNING ALGORITHMS
% =========================================================

\section{Machine Learning Algorithms}

% ---------------------------------------------------------
% LOGISTIC REGRESSION
% ---------------------------------------------------------

\subsection{Logistic Regression}

Logistic Regression is a probabilistic classification algorithm used
for binary classification.

The probability of the positive class is calculated using the sigmoid
function:

\[
P(y=1|x) =
\frac{1}{1+e^{-(w^Tx+b)}}
\]

where $w$ represents the model weights, $x$ represents the input
features and $b$ is the bias term.

A threshold is applied to convert the predicted probability into a
class label.

The model minimizes logistic or cross-entropy loss during training.

\subsubsection{Regularization}

Regularization is used to control model complexity and reduce
overfitting.

Two commonly used regularization techniques are:

\begin{itemize}
    \item \textbf{L1 Regularization:} Encourages some coefficients to
    become zero and can therefore perform feature selection.
    
    \item \textbf{L2 Regularization:} Penalizes large coefficients and
    generally provides smoother model weights.
\end{itemize}

The parameter $C$ controls the inverse strength of regularization.

A smaller value of $C$ results in stronger regularization, whereas a
larger value allows the model to fit the training data more closely.

% ---------------------------------------------------------
% SVM
% ---------------------------------------------------------

\subsection{Support Vector Machine}

Support Vector Machine is a margin-based classification algorithm.
It attempts to find an optimal decision boundary that separates the
classes while maximizing the margin between them.

Only a subset of training samples, known as support vectors,
directly determines the decision boundary.

The general SVM optimization objective can be represented as:

\[
\min_{w,b,\xi}
\frac{1}{2}\|w\|^2
+
C\sum_{i=1}^{n}\xi_i
\]

subject to:

\[
y_i(w^Tx_i+b)\geq1-\xi_i
\]

where $C$ controls the trade-off between maximizing the margin and
allowing classification errors.

% ---------------------------------------------------------
% SVM KERNELS
% ---------------------------------------------------------

\subsection{SVM Kernels}

Different kernels were evaluated in the experiment.

\begin{itemize}
    \item \textbf{Linear Kernel:} Suitable when the classes can be
    separated approximately using a linear decision boundary.
    
    \item \textbf{Polynomial Kernel:} Captures polynomial relationships
    between features.
    
    \item \textbf{RBF Kernel:} Can model complex non-linear decision
    boundaries.
    
    \item \textbf{Sigmoid Kernel:} Produces a decision function with
    behaviour similar to a sigmoid activation.
\end{itemize}

% =========================================================
% 7. EXPERIMENTAL SETUP
% =========================================================

\section{Experimental Setup}

\begin{longtable}{|p{6cm}|p{8cm}|}
\hline
\textbf{Parameter} & \textbf{Value} \\ \hline

Programming Language & Python \\ \hline

Machine Learning Library & Scikit-learn \\ \hline

Dataset & Spambase \\ \hline

Classification Type & Binary Classification \\ \hline

Validation Method & Five-Fold Cross-Validation \\ \hline

Models & Logistic Regression and SVM \\ \hline

\end{longtable}

% =========================================================
% 8. HYPERPARAMETER SELECTION
% =========================================================

\section{Hyperparameter Selection}

Hyperparameter tuning was performed to identify suitable configurations
for Logistic Regression and SVM.

% ---------------------------------------------------------
% LOGISTIC REGRESSION SEARCH
% ---------------------------------------------------------

\subsection{Logistic Regression Hyperparameter Tuning}

The parameters considered for Logistic Regression included:

\begin{itemize}
    \item Regularization: L1 and L2
    \item $C$: 0.01, 0.1, 1, 10 and 100
    \item Solver: liblinear and saga
\end{itemize}

The best configuration obtained was:

\begin{longtable}{|p{6cm}|p{8cm}|}
\hline
\textbf{Parameter} & \textbf{Value} \\ \hline

$C$ & 1 \\ \hline

Penalty & L2 \\ \hline

Solver & saga \\ \hline

Best CV Accuracy & 0.9242 \\ \hline

Tuning Time & 289.1025 seconds \\ \hline

\end{longtable}

\textbf{Observation:}

The best Logistic Regression configuration was obtained with
$C=1$, L2 regularization and the saga solver. The best cross-validation
accuracy during tuning was 0.9242. The tuning process required
approximately 289.10 seconds. The selected L2 regularization provides
a balance between model complexity and generalization.

% ---------------------------------------------------------
% SVM SEARCH
% ---------------------------------------------------------

\subsection{SVM Hyperparameter Tuning}

The SVM model was evaluated using linear, polynomial, RBF and sigmoid
kernels.

The best configuration obtained through hyperparameter tuning was:

\begin{longtable}{|p{6cm}|p{8cm}|}
\hline
\textbf{Parameter} & \textbf{Value} \\ \hline

$C$ & 10 \\ \hline

Gamma & scale \\ \hline

Kernel & RBF \\ \hline

Best CV Accuracy & 0.9332 \\ \hline

Tuning Time & 120.9516 seconds \\ \hline

\end{longtable}

\textbf{Observation:}

The best SVM configuration used the RBF kernel with $C=10$ and
gamma set to scale. This configuration achieved a cross-validation
accuracy of 0.9332. The tuning process required approximately
120.95 seconds. The RBF kernel provided better validation performance
during tuning than the other tested configurations.

% =========================================================
% 9. BASELINE LOGISTIC REGRESSION
% =========================================================

\section{Baseline Logistic Regression Results}

The baseline Logistic Regression model was evaluated on the test
dataset before hyperparameter tuning.

\begin{longtable}{|p{7cm}|c|}
\hline
\textbf{Metric} & \textbf{Value} \\ \hline

Accuracy & 0.9294 \\ \hline

Precision & 0.9209 \\ \hline

Recall & 0.8981 \\ \hline

F1 Score & 0.9093 \\ \hline

Training Time & 0.1428 seconds \\ \hline

\end{longtable}

\textbf{Observation:}

The baseline Logistic Regression model achieved an accuracy of
92.94\%. Precision was 92.09\%, while recall was 89.81\%.
The resulting F1-score was 90.93\%. The model required approximately
0.143 seconds for training, showing that Logistic Regression was
relatively efficient for this dataset.

% =========================================================
% 10. TUNED LOGISTIC REGRESSION
% =========================================================

\section{Tuned Logistic Regression Results}

After hyperparameter tuning, Logistic Regression achieved the
following results:

\begin{longtable}{|p{7cm}|c|}
\hline
\textbf{Metric} & \textbf{Value} \\ \hline

Accuracy & 0.9305 \\ \hline

Precision & 0.9211 \\ \hline

Recall & 0.9008 \\ \hline

F1 Score & 0.9109 \\ \hline

\end{longtable}

\textbf{Observation:}

Hyperparameter tuning produced a small improvement in Logistic
Regression performance. Accuracy increased from 92.94\% to 93.05\%.
Precision increased from 92.09\% to 92.11\%, while recall increased
from 89.81\% to 90.08\%. The F1-score also increased from 90.93\%
to 91.09\%. Therefore, the tuned model provided slightly better
overall performance than the baseline model.

% =========================================================
% 11. SVM KERNEL COMPARISON
% =========================================================

\section{SVM Kernel Comparison}

The performance of SVM was evaluated using four different kernels.

\begin{longtable}{|p{3cm}|c|c|c|c|c|}
\hline
\textbf{Kernel} &
\textbf{Accuracy} &
\textbf{Precision} &
\textbf{Recall} &
\textbf{F1 Score} &
\textbf{Training Time (s)}
\\ \hline

Linear &
0.9294 &
0.9209 &
0.8981 &
0.9093 &
0.5773
\\ \hline

Polynomial &
0.7796 &
0.9598 &
0.4601 &
0.6220 &
0.5557
\\ \hline

RBF &
0.9273 &
0.9277 &
0.8843 &
0.9055 &
0.2895
\\ \hline

Sigmoid &
0.8849 &
0.8599 &
0.8457 &
0.8528 &
0.3540
\\ \hline

\end{longtable}

\subsection{SVM Prediction Time}

\begin{longtable}{|p{4cm}|c|}
\hline
\textbf{Kernel} & \textbf{Prediction Time (s)} \\ \hline

Linear & 0.0253 \\ \hline

Polynomial & 0.0763 \\ \hline

RBF & 0.0864 \\ \hline

Sigmoid & 0.0573 \\ \hline

\end{longtable}

\textbf{Observation:}

The linear kernel achieved an accuracy of 92.94\% and an F1-score of
90.93\%. The RBF kernel produced a similar accuracy of 92.73\% and
an F1-score of 90.55\%, while requiring less training time than the
linear kernel.

The polynomial kernel obtained a relatively high precision of 95.98\%,
but its recall was only 46.01\%, resulting in a much lower F1-score of
62.20\%. This indicates that the polynomial kernel missed a significant
number of positive samples.

The sigmoid kernel achieved an accuracy of 88.49\% and an F1-score of
85.28\%. Overall, the linear kernel provided the strongest performance
among the untuned SVM kernels based on accuracy and F1-score.

% =========================================================
% 12. BEST SVM
% =========================================================

\section{Tuned SVM Results}

The best SVM configuration obtained through hyperparameter tuning was:

\[
C=10,\qquad
\gamma=\text{scale},\qquad
\text{kernel}=\text{RBF}
\]

The final test performance was:

\begin{longtable}{|p{7cm}|c|}
\hline
\textbf{Metric} & \textbf{Value} \\ \hline

Accuracy & 0.9207 \\ \hline

Precision & 0.9143 \\ \hline

Recall & 0.8815 \\ \hline

F1 Score & 0.8976 \\ \hline

Prediction Time & 0.0700 seconds \\ \hline

\end{longtable}

\textbf{Observation:}

The tuned SVM model used an RBF kernel with $C=10$ and gamma set to
scale. It achieved a test accuracy of 92.07\%, precision of 91.43\%,
recall of 88.15\% and an F1-score of 89.76\%. Although its cross-
validation accuracy was high, the final test accuracy was lower than
that of the tuned Logistic Regression model.

% =========================================================
% 13. LOGISTIC REGRESSION VS SVM
% =========================================================

\section{Logistic Regression versus SVM}

The final tuned models were compared using the test-set metrics.

\begin{longtable}{|p{5cm}|c|c|}
\hline
\textbf{Metric} &
\textbf{Logistic Regression} &
\textbf{SVM}
\\ \hline

Accuracy & 0.9305 & 0.9207 \\ \hline

Precision & 0.9211 & 0.9143 \\ \hline

Recall & 0.9008 & 0.8815 \\ \hline

F1 Score & 0.9109 & 0.8976 \\ \hline

\end{longtable}

\textbf{Observation:}

The tuned Logistic Regression model achieved higher test accuracy,
precision, recall and F1-score than the tuned SVM model. Logistic
Regression obtained an accuracy of 93.05\%, whereas SVM obtained
92.07\%.

The difference indicates that Logistic Regression generalized slightly
better on the test set for this experiment. However, SVM achieved a
higher cross-validation score during tuning, indicating that its
validation performance was strong even though its final test
performance was lower.

% =========================================================
% 14. FIVE-FOLD CROSS VALIDATION
% =========================================================

\section{Five-Fold Cross-Validation}

Five-fold cross-validation was performed to compare the consistency
of Logistic Regression and SVM.

In each fold, four parts of the dataset were used for training and the
remaining part was used for validation.

\begin{longtable}{|c|c|c|}
\hline
\textbf{Fold} &
\textbf{Logistic Regression} &
\textbf{SVM}
\\ \hline

1 & 0.9131 & 0.9305 \\ \hline

2 & 0.9326 & 0.9283 \\ \hline

3 & 0.9283 & 0.9413 \\ \hline

4 & 0.9239 & 0.9326 \\ \hline

5 & 0.9239 & 0.9359 \\ \hline

\textbf{Average} & \textbf{0.9244} & \textbf{0.9337} \\ \hline

\end{longtable}

\textbf{Observation:}

Logistic Regression achieved an average five-fold cross-validation
accuracy of 92.44\%, while SVM achieved a higher average accuracy of
93.37\%. SVM performed better than Logistic Regression in four of the
five folds, except Fold 2. The relatively small variation in the fold
accuracies indicates consistent performance across the validation
sets.

% =========================================================
% 15. CROSS VALIDATION ANALYSIS
% =========================================================

\section{Cross-Validation Analysis}

\begin{longtable}{|p{6cm}|c|c|}
\hline
\textbf{Model} &
\textbf{Average CV Accuracy} &
\textbf{Test Accuracy}
\\ \hline

Logistic Regression &
0.9244 &
0.9305
\\ \hline

SVM &
0.9337 &
0.9207
\\ \hline

\end{longtable}

\textbf{Observation:}

SVM obtained the higher average cross-validation accuracy of 93.37\%,
while Logistic Regression obtained 92.44\%. However, the final test
accuracy of Logistic Regression was higher at 93.05\%, compared with
92.07\% for SVM.

This difference shows that cross-validation performance and test-set
performance do not necessarily produce the same ranking. Based on the
final test results, Logistic Regression was the better-performing
model for this experiment.

% =========================================================
% 16. PERFORMANCE SUMMARY
% =========================================================

\section{Performance Summary}

\begin{longtable}{|p{5cm}|c|c|c|c|}
\hline
\textbf{Model} &
\textbf{Accuracy} &
\textbf{Precision} &
\textbf{Recall} &
\textbf{F1 Score}
\\ \hline

Baseline Logistic Regression &
0.9294 &
0.9209 &
0.8981 &
0.9093
\\ \hline

Tuned Logistic Regression &
\textbf{0.9305} &
\textbf{0.9211} &
\textbf{0.9008} &
\textbf{0.9109}
\\ \hline

Best SVM &
0.9207 &
0.9143 &
0.8815 &
0.8976
\\ \hline

\end{longtable}

\textbf{Observation:}

Hyperparameter tuning improved the Logistic Regression model slightly
across all four reported performance metrics. The tuned Logistic
Regression model achieved the highest test accuracy, precision, recall
and F1-score among the final models. Therefore, Logistic Regression was
selected as the best-performing model based on the final test-set
results.

% =========================================================
% 17. TRAINING AND TUNING TIME
% =========================================================

\section{Training and Tuning Time Analysis}

\begin{longtable}{|p{6cm}|c|}
\hline
\textbf{Operation} & \textbf{Time (seconds)} \\ \hline

Baseline Logistic Regression Training & 0.1428 \\ \hline

Logistic Regression Hyperparameter Tuning & 289.1025 \\ \hline

Linear SVM Training & 0.5773 \\ \hline

Polynomial SVM Training & 0.5557 \\ \hline

RBF SVM Training & 0.2895 \\ \hline

Sigmoid SVM Training & 0.3540 \\ \hline

SVM Hyperparameter Tuning & 120.9516 \\ \hline

\end{longtable}

\textbf{Observation:}

The baseline Logistic Regression model required only 0.1428 seconds
for training. Hyperparameter tuning required considerably more time
because multiple parameter combinations were evaluated using
cross-validation. Logistic Regression tuning required approximately
289.10 seconds, while SVM tuning required approximately 120.95
seconds. Among the SVM kernels, the RBF kernel had the lowest training
time at approximately 0.2895 seconds.

% =========================================================
% 18. DISCUSSION
% =========================================================

\section{Discussion}

The experiment compared a linear probabilistic classifier,
Logistic Regression, with a margin-based classifier, SVM.

The baseline Logistic Regression model already produced strong
performance with an accuracy of 92.94\%. After hyperparameter tuning,
its accuracy increased to 93.05\%. The improvement was small but was
also observed in precision, recall and F1-score.

Among the SVM kernels, the linear kernel achieved the highest accuracy
of 92.94\% before tuning. The RBF kernel achieved a comparable
accuracy of 92.73\% while requiring less training time. The polynomial
kernel produced high precision but considerably lower recall, which
resulted in a lower F1-score.

The tuned SVM used the RBF kernel with $C=10$ and gamma set to scale.
Its cross-validation accuracy was 93.32\%, but its final test accuracy
was 92.07\%.

The five-fold cross-validation results showed that SVM achieved a
higher average validation accuracy of 93.37\% compared with 92.44\%
for Logistic Regression. However, Logistic Regression achieved the
higher final test accuracy.

Therefore, based on the final test-set metrics, tuned Logistic
Regression was selected as the best-performing model.

% =========================================================
% 19. STRENGTHS AND LIMITATIONS
% =========================================================

\section{Strengths and Limitations}

\subsection{Logistic Regression}

\textbf{Strengths:}

\begin{itemize}
    \item Simple and computationally efficient.
    \item Easy to interpret.
    \item Performs well for linearly separable classification problems.
    \item Regularization helps control model complexity.
\end{itemize}

\textbf{Limitations:}

\begin{itemize}
    \item Primarily provides a linear decision boundary.
    \item Performance may decrease when complex non-linear relationships
    exist between features.
\end{itemize}

\subsection{SVM}

\textbf{Strengths:}

\begin{itemize}
    \item Effective for binary classification.
    \item Can model non-linear relationships through kernels.
    \item Maximizes the margin between classes.
\end{itemize}

\textbf{Limitations:}

\begin{itemize}
    \item Hyperparameter tuning can be computationally expensive.
    \item Kernel selection affects performance significantly.
    \item The resulting model is less directly interpretable than
    Logistic Regression.
\end{itemize}

% =================================================
% SVM KERNEL-WISE ACCURACY
% =================================================

\subsection{SVM Kernel-wise Accuracy}

\begin{figure}[H]
    \centering
    \includegraphics[width=0.75\linewidth]{04_svm_kernel_accuracy.png}
    \caption{Accuracy comparison of different SVM kernels}
    \label{fig:svm_kernel_accuracy}
\end{figure}

\textbf{Observation:}

The linear kernel achieved the highest accuracy of 92.94\%,
followed closely by the RBF kernel with an accuracy of 92.73\%.
The sigmoid kernel achieved an accuracy of 88.49\%, while the
polynomial kernel produced the lowest accuracy of 77.96\%.
Therefore, the linear kernel provided the best classification
accuracy among the four SVM kernels evaluated.


% =================================================
% SVM KERNEL-WISE F1 SCORE
% =================================================

\subsection{SVM Kernel-wise F1 Score}

\begin{figure}[H]
    \centering
    \includegraphics[width=0.75\linewidth]{05_svm_kernel_f1.png}
    \caption{F1-score comparison of different SVM kernels}
    \label{fig:svm_kernel_f1}
\end{figure}

\textbf{Observation:}

The linear kernel achieved the highest F1-score of 0.9093,
followed closely by the RBF kernel with an F1-score of 0.9055.
The sigmoid kernel obtained an F1-score of 0.8528, whereas
the polynomial kernel obtained the lowest F1-score of 0.6220.
The low F1-score of the polynomial kernel is mainly associated
with its low recall value of 0.4601.

% =================================================
% LOGISTIC REGRESSION VS SVM
% =================================================

\subsection{Logistic Regression versus SVM}

\begin{figure}[H]
    \centering
    \includegraphics[width=0.85\linewidth]{06_model_comparison.png}
    \caption{Performance comparison of Logistic Regression and SVM}
    \label{fig:logistic_vs_svm}
\end{figure}

\textbf{Observation:}

The tuned Logistic Regression model achieved an accuracy of
93.05\%, whereas the best SVM model achieved an accuracy of
92.07\%. Logistic Regression also obtained higher recall and
F1-score than SVM. The precision values of both models were
relatively close. Based on the final test-set results, Logistic
Regression provided better overall performance than SVM.

% =================================================
% FIVE-FOLD CROSS-VALIDATION
% =================================================

\subsection{Five-Fold Cross-Validation}

\begin{figure}[H]
    \centering
    \includegraphics[width=0.80\linewidth]{07_cross_validation.png}
    \caption{Five-fold cross-validation accuracy comparison}
    \label{fig:five_fold_cv}
\end{figure}

\textbf{Observation:}

The five-fold cross-validation results show that SVM achieved
a higher average accuracy of 93.37\% compared with 92.44\%
for Logistic Regression. SVM obtained its highest accuracy
of 94.13\% in Fold 3, while Logistic Regression obtained its
highest accuracy of 93.26\% in Fold 2. Overall, SVM showed
slightly better average validation performance across the
five folds.

% =========================================================
% 20. FINAL OBSERVATION
% =========================================================

\section{Final Observation}

The experiment demonstrated the application of Logistic Regression and
Support Vector Machine for binary spam and ham classification.

The tuned Logistic Regression model achieved a test accuracy of
\textbf{93.05\%}, precision of \textbf{92.11\%}, recall of
\textbf{90.08\%}, and F1-score of \textbf{91.09\%}. Its best
hyperparameters were $C=1$, L2 regularization and the saga solver.

The best SVM configuration used an RBF kernel with $C=10$ and gamma
set to scale. It achieved a test accuracy of \textbf{92.07\%},
precision of \textbf{91.43\%}, recall of \textbf{88.15\%}, and
F1-score of \textbf{89.76\%}.

Five-fold cross-validation resulted in average accuracies of
\textbf{92.44\%} for Logistic Regression and \textbf{93.37\%} for
SVM. Although SVM obtained a higher average cross-validation accuracy,
Logistic Regression achieved better performance on the final test set.

Hence, based on the final test-set performance, \textbf{tuned Logistic
Regression was selected as the best-performing classifier for this
experiment}.

% =========================================================
% 21. CONCLUSION
% =========================================================

\section{Conclusion}

Binary classification was successfully performed using Logistic
Regression and Support Vector Machine on the Spambase dataset.

Different SVM kernels were evaluated and hyperparameter tuning was
performed for both models. Hyperparameter tuning improved the
performance of Logistic Regression slightly. SVM showed strong
cross-validation performance, particularly with the RBF kernel.

The final comparison showed that tuned Logistic Regression achieved
the highest test accuracy of 93.05\% and the highest F1-score of
91.09\%. Therefore, Logistic Regression provided the best final
test-set performance in this experiment.

The experiment also demonstrated the importance of hyperparameter
selection, kernel selection and cross-validation when evaluating
machine learning classification models.

% =========================================================
% 22. LEARNING OUTCOMES
% =========================================================

\section{Learning Outcomes}

After completing this experiment, the following concepts were
understood:

\begin{itemize}
    \item Binary classification using Logistic Regression.
    \item Margin-based classification using SVM.
    \item Effect of L1 and L2 regularization.
    \item Effect of the $C$ hyperparameter.
    \item Behaviour of different SVM kernels.
    \item Hyperparameter tuning using cross-validation.
    \item Evaluation using accuracy, precision, recall and F1-score.
    \item Five-fold cross-validation.
    \item Comparison of validation performance and test performance.
\end{itemize}

% =========================================================
% 23. REFERENCES
% =========================================================

\section{References}

\begin{enumerate}
    \item Scikit-learn Documentation, Logistic Regression.
    
    \item Scikit-learn Documentation, Support Vector Machines.
    
    \item Scikit-learn Documentation, Hyperparameter Optimization.
    
    \item Spambase Dataset, UCI Machine Learning Repository.
    
    \item Spambase Dataset, Kaggle.
    
    \item ICS1512 Machine Learning Algorithms Laboratory,
    Experiment 4 Manual.
\end{enumerate}

\end{document}